\section{Weil's full criterion}
\label{sec:weil-full}

A positive realization of the zeros must recover a full arithmetic
quadratic form. The criterion used below is the classical explicit
formula and positivity theorem of Weil; its general contour and
zero-separation theorem are imported in the stated test class. The
normalization calculation records every local term. The subsequent
narrow-window estimates are direct inequalities for this full form.

\begin{definition}[Weil test class]
\label{def:weil-class}
Let \(\mathcal W\) consist of complex-valued functions
\(f:\mathbb R\to\mathbb C\) which are continuously differentiable away
from finitely many points, have finite one-sided limits there for both
\(f\) and \(f'\), and satisfy the convention
\[
 f(a)=\tfrac12\bigl(f(a+)+f(a-)\bigr)
\]
at every exceptional point \(a\).  Suppose also that, for some \(b>0\),
\[
 f(u),f'(u)=O(e^{-(1/2+b)|u|}).
\]
These hypotheses are Weil's conditions (A)--(B)
\cite[p.~259]{Weil1952}.
For \(f\in\mathcal W\), set
\[
 \widetilde f(u)=\overline{f(-u)},\qquad g=f*\widetilde f,
\]
and fix the Fourier convention
\[
 h(t)=\widehat g(t)=\int_{\mathbb R}g(u)e^{itu}\,du,
 \qquad
 g(u)=\frac1{2\pi}\int_{\mathbb R}h(t)e^{-itu}\,dt.
\]
\end{definition}

For every \(0<b'<b\), the convolution \(g\) satisfies the same conditions
with \(b'\) in place of \(b\).  Piecewise integration by parts, uniformly
for \(|\eta|\leq\tfrac12\), gives
\[
 \widehat f(t+i\eta)=O((1+|t|)^{-1}).
\]
The Fourier transform extends to the strip needed below, and
\[
 h(z)=\widehat f(z)\overline{\widehat f(\overline z)},
 \qquad h(t)=|\widehat f(t)|^2\quad(t\in\mathbb R).
\]
Equivalently, for the centered Mellin transform
\[
 \Phi_f(s)=\int_{\mathbb R}f(u)e^{(s-1/2)u}\,du,
\]
one has
\[
 h\!\left(-i(s-\tfrac12)\right)
 =\Phi_f(s)\overline{\Phi_f(1-\overline s)}.
\]

Write $\psiGamma=\Gamma'/\Gamma$ for the digamma function.
The von Mangoldt function is $\vmLambda(n)=\log p$ when $n=p^a$
with $a\ge1$, and zero otherwise.

\begin{theorem}[Weil explicit formula and positivity criterion]
\label{thm:weil}
Let \(f\in\mathcal W\), \(g=f*\widetilde f\), and
\(h=\widehat g\).  Write a nontrivial zero as
\(\rho=\beta_\rho+i\gamma_\rho\), with multiplicity \(m(\rho)\), and set
\[
 \sum_\rho^{*}h\!\left(-i(\rho-\tfrac12)\right)
 :=\lim_{\substack{T\to\infty\\
                    T\notin\{|\gamma_\rho|\}}}
 \sum_{|\gamma_\rho|<T}m(\rho)
 h\!\left(-i(\rho-\tfrac12)\right).
\]
This limit is in fact absolute for the stated convolution-square class.
Then
\[
\begin{aligned}
 \sum_\rho^{*}h\!\left(-i(\rho-\tfrac12)\right)
 ={}&h(i/2)+h(-i/2)\\
 &+\frac1{2\pi}\int_{\mathbb R}h(t)
 \left[\operatorname{Re}\psiGamma\!\left(\frac14+\frac{it}{2}\right)
 -\log\pi\right]dt\\
 &-\sum_{n\geq2}\frac{\vmLambda(n)}{\sqrt n}
 \bigl(g(\log n)+g(-\log n)\bigr),
\end{aligned}
\]
The Riemann Hypothesis
holds if and only if this full common value is nonnegative for every
\(f\in\mathcal W\).
\end{theorem}

\begin{proof}[Normalization and consequences of the classical theorem]
Moving the logarithmic-derivative contour across the two poles of
\(\Lambda(s)\) gives \(h(i/2)+h(-i/2)\).  Fourier inversion applied to
\[
 -\frac{\zeta'}{\zeta}(s)
 =\sum_{n\geq2}\frac{\vmLambda(n)}{n^s}
\]
gives the two prime-power terms.  Finally,
\[
 \frac{\GammaR'}{\GammaR}(s)
 =-\frac12\log\pi+\frac12\psiGamma(s/2),
\]
so the two sides of the critical line give the kernel
\[
 2\operatorname{Re}\frac{\GammaR'}{\GammaR}
 \left(\frac12+it\right)
 =\operatorname{Re}\psiGamma\!\left(\frac14+\frac{it}{2}\right)-\log\pi.
\]
The decay in Definition~\ref{def:weil-class} justifies the contour passage,
the prime sum, and the archimedean integral.  The estimate preceding the
theorem gives
\(h(-i(\rho-\tfrac12))=O((1+|\gamma_\rho|)^{-2})\); together with
the standard bound \(N(T)=O(T\log T)\), it also proves the asserted absolute
convergence.  This is Weil's formula~(11)
in the present Fourier normalization
\cite[p.~262, formula~(11)]{Weil1952}.

If the Riemann Hypothesis holds, write
\(\rho=\tfrac12+i\gamma\); every zero contributes
\(h(\gamma)=|\widehat f(\gamma)|^2\geq0\).  Conversely, Weil's
Gaussian-polynomial construction associated with an off-line zero produces
an \(f\in\mathcal W\) for which the zero-side sum is negative
\cite[pp.~262--263, unnumbered lemma immediately following
formula~(11)]{Weil1952}.
\end{proof}

When \(g\) is real and even, the last line of the formula becomes
\[
 -2\sum_p\sum_{m\geq1}\frac{\log p}{p^{m/2}}g(m\log p).
\]
Both the polar terms and this prime sum are essential to the criterion.

\begin{proposition}[The gamma term is not positive by itself]
\label{prop:gamma-negative}
Set
\[
 k(t)=\operatorname{Re}\psiGamma\!\left(\frac14+\frac{it}{2}\right)-\log\pi.
\]
There are functions \(f\in C_c^\infty(\mathbb R)\subset\mathcal W\) for
which
\[
 \frac1{2\pi}\int_{\mathbb R}|\widehat f(t)|^2k(t)\,dt<0.
\]
\end{proposition}

\begin{proof}
The special value
\[
 k(0)=-\gamma-\frac\pi2-3\log2-\log\pi
\]
is negative.  Choose nonzero \(\varphi\in C_c^\infty(\mathbb R)\) and put
\(f_R(u)=R^{-1/2}\varphi(u/R)\).  Then
\[
 |\widehat f_R(t)|^2=R|\widehat\varphi(Rt)|^2
\]
and therefore
\[
 \frac1{2\pi}\int_{\mathbb R}|\widehat f_R(t)|^2k(t)\,dt
 =\frac1{2\pi}\int_{\mathbb R}|\widehat\varphi(v)|^2k(v/R)\,dv.
\]
Since \(k(t)=O(\log(2+|t|))\) and \(\widehat\varphi\) is rapidly
decreasing, dominated convergence makes the last expression tend to
\[
 \frac{k(0)}{2\pi}
 \int_{\mathbb R}|\widehat\varphi(v)|^2\,dv<0.
\]
It is therefore negative for all sufficiently large \(R\).
\end{proof}

Proposition~\ref{prop:gamma-negative} shows that the isolated gamma term is
not nonnegative on all Weil convolution squares and therefore cannot replace
the full form in Theorem~\ref{thm:weil}.  The theorem-level criterion uses
the entire test class and every local term.
